\documentclass[11pt]{article}
\usepackage[margin=0.85in]{geometry}
\usepackage{amsmath,amssymb,booktabs,longtable,array,xurl,hyperref}
\hypersetup{hidelinks}
\setlength{\parindent}{0pt}
\setlength{\parskip}{6pt}
\title{Joint Dynamic Output-Feedback Pipeline}
\author{Pan Dynamics Collision Avoidance Research}
\date{November 4, 2026}
\begin{document}
\maketitle

\section*{Implemented behavior}
The primary campaign entry points now run raw-sensor dynamic output feedback.
The controller conditions joint physical hypotheses on complete GNSS-position,
body-acceleration, gyro and radar blocks, then optimizes both actuator and
memory maps:
\[
 u_i=v_i+K_i m_i,\qquad m_{i+1}=\tanh(A_i m_i+B_i\varphi(Y_{i+1})+c_i).
\]
Every scenario shares the stage coefficients and the current known memory.
The selected first memory map is stored with its command and applied exactly
once to the next complete sensor block. Persistent tire/friction scales and
sensor biases remain attached to their physical hypotheses. The primary path
calls no NRMM runtime, observer gain synthesis, force inversion or CLF table.
The state-input API and separate observers remain explicit research comparators.

The actual implementation eliminates shooting states by nonlinear numerical
rollout. All finite scenarios receive hard actuator, command-increment,
physical-domain, rectangle-separation and terminal constraints. The last stage
is fixed to the physical fading-brake backup. A regularized Gauss--Newton QP
proposes one online SQP update. Nonlinear reevaluation and backtracking govern
acceptance. Initial feasibility search may use elastic QP rows, but a command
is committed only after the original finite hard constraints pass. A shifted
plan is rechecked after conditioning; it is never presumed feasible from a
newly fitted or changed information representation.

\section*{Theory and physical-model corrections}
The rewritten \path{controller/CONTROLLER_THEORY.md} states the exact joint
transition, complete-block posterior, hard predictive problem, and conditional
recursive-feasibility proof. That proof requires sound uncertainty coverage,
whole-hold verification, nested posterior containment, and policy/set/certificate
classes capable of retaining the shift. Finite scenarios do not meet these
universal requirements.

The terminal component uses $u_f=(0,-b_0)$ and, on the pure low-speed branch,
\[
 s_E=|v_x|+\tfrac{\tau_\ell}{2}|v_y r|,\qquad
 D_E=s_E/\lambda+\tau_\ell|v_y|.
\]
Its dissipation bounds remaining reference-point travel. The forward-domain
restrictions $v_x\ge0$ and $v_y r\ge0$ prevent reverse forcing at zero speed.
A pose-preserving control-boundary rest reset reduces $D_E$ and preserves the
separation reserve. Sensor-only boundaries do not apply that reset.
For persistent constant target acceleration $A_T<0$, the sharper travel term
$D_T=V_T^2/(-2A_T)$ replaces the unnecessary global lower-bound denominator;
absorbed targets have $D_T=0$ and cannot restart from latent acceleration.

The terminal component does not support continuing targets, restart,
persistent residual forces, multiple targets, or hard road constraints.
The smoothed absolute-value entry inequalities are sufficient for entry to the
exact terminal set, but their own representation is not asserted invariant.
The exact theorem does not automatically transfer to this smoothed, sampled,
finite-scenario implementation.

\section*{Validation results}
MATLAB R2026a Update 3 and Optimization Toolbox were used. The focused suite
passed all 223 tests, with zero failures and zero incomplete tests. It covers
the new raw-block behaviors and existing configuration, source-budget,
Fiala-tire, road-load, RTI-model, terminal-set and input-geometry behavior.
Twenty new behavior tests include shared root actions, observation-dependent
future commands, exact committed memory, full-block consistency despite equal
compressed features, previous-input acceleration, dropout retention, persistent
parameters/bias, absorbing target stop, control-only rest reset, empty-information
failure, unverified-plan rejection, revalidated fallback, and the default
campaign routing. An independent tight-tolerance \texttt{ode45} calculation
checks the low-speed travel and separation-reserve inequalities.

The closed-loop runs used seed 20261104. Sensor noise was bounded, sampled
within half of the declared sensor radii, with no GNSS velocity measurement.
The three-member initial experimental information contract includes the actual
initial member by construction. It is not an enclosure inferred from raw data.
Subsequent truth states are used only for sensor synthesis and offline replay.
Forty-one replay instants per hold provide empirical clearance checks.

\begin{center}\small
\begin{tabular}{@{}lrrrr@{}}
\toprule
Case & Speed & Holds & Min. gap & Median call \\
 & m/s & executed & m & s \\
\midrule
Stopped lead & 1.5 & 8 & 28.7431 & 1.642 \\
Braking lead & 1.5 & 8 & 25.2000 & 1.552 \\
Stopped lead, dropout & 1.5 & 6 & 28.8311 & 1.550 \\
Stopped lead, moving ego & 8 & 3 & 26.9349 & 9.515 \\
Stopped lead, admission & 25 & 0 & --- & --- \\
Braking lead, admission & 25 & 0 & --- & --- \\
\bottomrule
\end{tabular}
\end{center}

The 1.5 m/s runs use four 0.05 s holds, two memory coordinates, three current
hypotheses and three fixed future noise directions. Dropouts occur on holds
2 and 3. The 8 m/s run uses a 0.1 s hold, 0.025 s sensor period, one future noise
direction, and 17 stages selected once at initialization. All accepted recorded
plans have zero maximum positive inequality violation in the stored numerical
evaluation. Maximum measured calls are 2.580 s for the low-speed stopped lead,
2.247 s for the low-speed braking lead, 2.065 s for dropout and 10.001 s for the
8 m/s run. These are short functional checks, not encounter completion,
cruise recovery, or robust-safety evidence.

Both 25 m/s primary challenge cases failed initial admission with a five-second
search budget. The reported maximum remaining hard-constraint violations are
2.56521 and 2.59021 respectively; neither case issued an actuator command.
They establish failure of this initialization/search configuration, not physical
infeasibility. The finite-difference dense SQP solver is substantially slower
than the control hold. Budget checks between operations cannot interrupt a
running QP. No real-time-performance claim is made.

\section*{Reproduction and artifacts}
Run from the repository root. Raw results and settings are retained under
\path{report/UNIFIED_OUTPUT_FEEDBACK_20261104/}. Each run stores its configuration,
explicit initial hypotheses, seed, raw blocks, applied commands, memory, survivor
counts, constraint checks, timing and offline clearance. Successful runs are
\path{low-speed/}, \path{dropout/}, and \path{moving-ego/}; the failed initial
admissions are in \path{high-speed-admission/}. The complete focused test list
and actual outcomes are in \path{tests.csv}. Exact executed commands are in
\path{commands.txt}.

The principal source changes are in \path{controller/solvePredictiveControl.m},
\path{controller/collisionAvoidanceController.m},
\path{controller/nonlinearBicycleModel.m},
\path{controller/terminalSafeSet.m}, and
\path{config/collisionAvoidanceControllerConfig.m}.
The raw interface and limits are documented in \path{README.md}; the independent
plant/sensor driver is \path{scripts/runJointOutputFeedbackScenario.m}.

The earlier probe's missing CLF-table dependency was removed from joint
initialization: its geometric seed now needs only route/speed data. The first
braking-target probe exposed an excessively loose global target-deceleration
floor; using the persistent acceleration coordinate corrected that bound.
Neither preliminary failure is represented as a successful run.

\section*{Limits and next engineering work}
Finite hypotheses can lose all members even when an unrepresented physical
state satisfies the bounded-noise contract. Fixed future noise directions do
not cover all measurement histories. The RK4 flow and sampled speed/rate checks
are not validated continuous enclosures. The intersample reserve is conditional
on valid continuous speed/rate caps. These limitations remain after nonlinear
finite-scenario acceptance. No deterministic or probabilistic collision-safety
claim follows from this implementation.

The next work is to improve feasible initialization for the 25 m/s steering
challenge, replace dense finite differences with sparse joint sensitivities,
and add a sound information representation and verifier if a universal
certificate is required. Continuing-target and road-safe terminal components
need separate proofs. The mathematical formulation and implementation do not
claim global optimality, observer convergence, or tracking convergence.
\end{document}
